\section{The constrained finite-dimensional metric evolution}\label{sec:finite-flow}

Let $\cP_r=\{H=H^*>0:\det H=1\}$. We use its affine invariant metric
and distance
\begin{equation}\label{eq:matrix-metric}
 g_H(V,W)=\tr(H^{-1}VH^{-1}W),\qquad
 d_r(H,K)^2=\tr\!\left(\log(H^{-1/2}KH^{-1/2})^2\right).
\end{equation}
Its geodesic is
\[
 \Geo(H,K;u)=H^{1/2}(H^{-1/2}KH^{-1/2})^uH^{1/2}.
\]
This is a complete simply connected space of nonpositive curvature;
the distance, geodesic and exponential metric inequality are reviewed
in \cite[Section 2.3]{bhatia}.
Congruence by an invertible matrix preserves its metric whenever the
congruence preserves the determinant-one locus. In particular unitary
transport is an isometry. Logarithms such as $\log(H^{-1}K)$ are
self-adjoint in the $H$ metric; they need not be Hermitian in the
identity frame.

The configuration space dictated by the bundle lift is
\begin{equation}\label{eq:finite-space}
 \cF=\{(H_0,H_1,H_2,H_3)\in\cP_2^4:
                   H_{-j}=JH_jJ^*\}.
\end{equation}
It is the fixed set of an isometric involution of the product with metric
$\sum_jA_jg_{H_j}$. It is therefore totally geodesic and geodesically
convex. More explicitly it is
$\cP_2^J\times\cP_2\times\cP_2^J$, with middle weight $2A_1$.
Each endpoint factor consists of $\exp(ae_x)$, $a\in\RR$, so
$\dim_\RR\cF=5$. The set is independent of $\theta$ in these frames;
the energy contains the parameter.

Put
\begin{equation}\label{eq:finite-potential}
 \Phi_{t,\theta}(H)=\frac\pi L\sum_{j\bmod4}
                d_2(H_{j+1},RH_jR^*)^2.
\end{equation}
Every edge appears once. The relation $JR=R^{-1}J$ makes this function
invariant under the involution defining \eqref{eq:finite-space}.

\begin{proposition}\label{prop:finite-flow}
The negative gradient flow of \eqref{eq:finite-potential} on $\cF$ is
the restriction of the four-vertex system
\begin{equation}\label{eq:finite-flow}
 A_jH_j^{-1}\dot H_j=\frac{2\pi}{L}
 \left\{\log(H_j^{-1}RH_{j-1}R^*)
             +\log(H_j^{-1}R^*H_{j+1}R)\right\}.
\end{equation}
For initial tuples in a fixed bounded set, it has a solution for every
physical time $s\ge0$, confined to a compact set independent of $t$ and
$\theta$. Distances between two solutions with the same parameters are
nonincreasing. The speed is $O(L^{-1})$, uniformly for all $s\ge0$;
the same bound holds for the time derivative of any fixed smooth tuple
function on the confined range.
\end{proposition}

\begin{proof}
The negative gradient in the first variable of
$\tfrac12d_2(H,K)^2$ is $H\log(H^{-1}K)$. Differentiating the two
incident terms and dividing by $A_j$ gives \eqref{eq:finite-flow}.
An invariant function has tangent gradient on the fixed set of an
isometric involution, so \eqref{eq:finite-space} is preserved. The
displayed logarithms are trace-free; the determinant condition is preserved.

Squared distance is jointly geodesically convex in this nonpositively
curved matrix space. Hence $\Phi$ is convex, and the identity tuple $I$
is a minimizer. The first-variation inequality for a convex function
gives
\[
 \frac{d}{ds}\frac12d_M(H(s),I)^2\le-\Phi(H(s))\le0,
\]
where $d_M^2=\sum_jA_jd_2^2$ is the product squared distance.
Closed bounded sets are compact in this finite-dimensional space,
proving confinement and continuation for all time. Applying the same
convexity inequality to the two geodesics between two solutions gives
nonexpansion. On the confined range every logarithm in
\eqref{eq:finite-flow} is bounded uniformly in the unitary $R$.
The equation then gives the speed and derivative bounds.
\end{proof}

The system is genuinely noncommuting. For example, choose $ab\ne0$ and
\[
 H_0=e^{a\sigma_x},\qquad H_1=e^{b\sigma_z},\qquad
 H_2=I,\qquad H_3=Je^{b\sigma_z}J^*=e^{-b\sigma_z}.
\]
This lies in $\cF$. In the common frame of the first edge,
$H_0$ and $R^*H_1R=H_1$ do not commute. Simultaneous frame conjugation
cannot remove this commutator, so the construction is not restricted
to independent scalar equations.

\begin{lemma}\label{lem:finite-decay}
Set $Z_j=\log H_j$,
$R_f=\tfrac12\sum_jA_j\|Z_j\|_{\HS}^2$, and
$\nu_0=(2\pi/L)\mu_0(\theta)$. Along \eqref{eq:finite-flow},
\begin{equation}\label{eq:finite-decay}
 \Phi(H)\ge\nu_0R_f,\qquad
 R_f'\le-2\Phi(H)\le-2\nu_0R_f.
\end{equation}
The linearization at the identity is the negative of the effective
operator with energy $(2\pi/L)q_\theta$ and mass $m_A$.
\end{lemma}

\begin{proof}
The logarithm based at the identity does not increase distance, so
\[
 d_2(H_{j+1},RH_jR^*)^2
       \ge\|Z_{j+1}-\Ad(R)Z_j\|^2.
\]
The logarithmic tuple satisfies \eqref{eq:linear-invariance}.
Using \eqref{eq:five-forms} gives the first inequality.

Radial contraction $H_j\mapsto e^{uZ_j}$ stays in $\cF$ and satisfies
$\Phi(e^{uZ})\le u^2\Phi(H)$ for $0\le u\le1$. This is the
distance contraction toward the common fixed point $I$. Differentiation
from the left at $u=1$ gives a radial derivative at least $2\Phi(H)$.
Pairing the negative gradient with the gradient of $R_f$ proves
\eqref{eq:finite-decay}. Finally expansion of the logarithms in
\eqref{eq:finite-flow} at $I$ gives the transported connection-graph
operator, whose restriction was computed in Section~\ref{sec:spectrum}.
\end{proof}

The parameter $\nu_0$ includes the exact zero mode at reducibility.
Thus \eqref{eq:finite-decay} gives a uniform exponential tail in $s/L$
on classes $\mu_0(\theta)\ge\kappa>0$. It does not give such a tail
when $\mu_0\to0$, although confinement and nonexpansion remain uniform.
